\section{Rigidity of each nonzero block}
\label{sec:blocks}

A periodic orbit of $F$ is encoded by a bi-infinite periodic integer
sequence satisfying
\begin{equation}
 x_{i+n}+x_i=\prod_{j=1}^{n-1}x_{i+j}+a.
 \label{eq:recurrence}
\end{equation}
Its consecutive length-$n$ blocks are the native states. In this
section $n\ge4$, $A=|a|$, and $B=A+4$.

\begin{lemma}[Internal maxima]
\label{lem:internal}
If a periodic sequence has no zero, all its coordinates have absolute
value at most $B$. If it has zeros, let
$x_1,\ldots,x_r$ be a maximal nonzero block, with
$x_0=x_{r+1}=0$, and put $M=\max_{1\le j\le r}|x_j|$.
When $r\ge n-1$ and $M>B$, every occurrence of absolute value $M$
in this block is an endpoint.
\end{lemma}
\begin{proof}
For a nonzero window of length $n-1$ contained in this block,
its two external neighbours belong either to the block or to its
adjacent zeros. Therefore
\begin{equation}
 \left|\prod_{j=1}^{n-1}x_{i+j}\right|\le2M+A.
 \label{eq:local-product-bound}
\end{equation}
This is a block-local estimate, independent of the heights in other
blocks beyond the zeros.

Suppose an internal position has absolute value $M>A$. Because
$n-1\ge3$, choose a length-$(n-1)$ window within the block which
contains that position strictly internally. Write its external
neighbours as $z_-,z_+$. Moving the window one place left preserves
the maximum factor and all nonzero internal factors. If the new
factor $z_-$ is zero, the inequality $M|z_-|\le2M+A$ is immediate;
otherwise the shifted window is in the same block and
\eqref{eq:local-product-bound} proves it. The right shift similarly
gives $M|z_+|\le2M+A$.
Since $M>A$, both neighbours have absolute value less than three,
and thus at most two. The unshifted recurrence now gives
\[
 M\le\left|\prod x_j\right|=|z_-+z_+-a|\le A+4=B.
\]
This contradicts $M>B$. If the whole sequence is nonzero, use its
global maximum in an internal window; the identical argument proves
the first assertion. The all-zero sequence needs no block analysis.
\end{proof}

\begin{lemma}[Large blocks]
\label{lem:large-block}
Every maximal nonzero block of length at least $n-1$ containing a
coordinate of absolute value greater than $B$ has length exactly $n$.
Its $n-2$ internal coordinates are units. After possibly reversing
indices in the proof, its endpoints satisfy
\begin{equation}
 x_n=\varepsilon x_1+a,\qquad
 \varepsilon=\prod_{j=2}^{n-1}x_j\in\{1,-1\},\qquad
 (1+\varepsilon)a=0.
 \label{eq:block-relations}
\end{equation}
\end{lemma}
\begin{proof}
If the block length is $n-1$, the recurrence across its two zeros
gives $\prod_{j=1}^{n-1}x_j=-a$. The product is a nonzero integer,
so $a\ne0$ and every factor has absolute value at most $A$.
It cannot be a large block.

Suppose the length is at least $n$ and its maximum is $M>B$.
Lemma~\ref{lem:internal} places a maximum at an endpoint.
The recurrence is preserved by index reversal with the same $a$,
so assume $|x_1|=M$. This reversal is a proof device, not an
identification of oriented native orbits. Put
$C=\prod_{j=2}^{n-1}x_j$. The recurrence at the left zero gives
$x_n=Cx_1+a$, whence
\[
 M|C|\le M+A<2M.
\]
The nonzero integer $C$ is therefore a sign $\varepsilon$, and each
of its factors is a unit. If the length equals $n$, the next
recurrence gives $x_1=\varepsilon x_n+a$, proving
\eqref{eq:block-relations}. The second endpoint need not also
exceed $B$; no such assumption is used.

It remains to exclude a length at least $n+1$. The next coordinate is
\[
 x_{n+1}=\varepsilon x_n+a-x_1=(1+\varepsilon)a.
\]
It is nonzero, so $a\ne0$, $\varepsilon=1$, and $x_{n+1}=2a$.
The window $x_3,\ldots,x_{n+1}$ has unit factors up to $x_{n-1}$
and remaining factors $x_n,2a$. Its external neighbours are the
unit $x_2$ and a coordinate $x_{n+2}$ of absolute value at most $M$
(possibly the right zero). Consequently
\[
 2A|x_n|\le M+A+1,\qquad
 (2A-1)M\le2A^2+A+1.
\]
Here $|x_n|=|x_1+a|\ge M-A$ and $A\ge1$. But
\[
 (2A-1)(A+4)-(2A^2+A+1)=6A-5>0,
\]
contradicting $M>A+4$. A right-end maximum is covered by the same
reversal argument.
\end{proof}

\begin{corollary}[Every update is linear after finite tagging]
\label{cor:window}
In any periodic sequence, each nonzero update product in
\eqref{eq:recurrence} either has all factors of absolute value at
most $B$, or has exactly one factor exceeding $B$ in absolute value,
with all its other factors in $\{1,-1\}$.
\end{corollary}
\begin{proof}
The zero-free case follows from Lemma~\ref{lem:internal}. Otherwise
a nonzero window lies in one maximal nonzero block. Blocks shorter
than $n-1$ cannot contain such a window; blocks of length $n-1$ are
bounded. A block with a large factor has exactly $n$ entries by
Lemma~\ref{lem:large-block}. Its two possible nonzero windows of
length $n-1$ each omit one endpoint. Thus a window with a large
factor has only one large endpoint, and all companions are units.
If neither endpoint in that window is large, it is the first case.
\end{proof}

The corollary includes arbitrarily many large coordinates separated
by zeros: an update product containing a zero is zero regardless of
the other factors. It imposes no independence assumption on distinct
blocks, no bound on their number, and no guess about the dimension
of a channel. These are the cases which a zero-free height bound
alone leaves untreated.
